\section{Infinite-dimensional examples and spherical generators}
\label{sec:E}

\subsection{The six-element inflation: a complete homotopy identification}
All products in the following statement are monoid products, with $gf$ denoting
the composite of the arrow $f$ followed by the arrow $g$. Consequently the
normalized nerve boundary on a chronological word $(m_1,\ldots,m_n)$ composes
the adjacent entries as $m_{i+1}m_i$.

\begin{theorem}\label{thm:E-wedge}
Let $M=\{1,a,b,c,d,y\}$, where $a=x_{00}$, $b=x_{01}$,
$c=x_{10}$ and $d=x_{11}$. The identity is $1$, and
\[
 x_{ij}x_{k\ell}=x_{i\ell},\qquad
 yx_{ij}=x_{0j},\qquad x_{ij}y=x_{i0},\qquad y^2=a.
\]
Then
\[
 BM\simeq\bigvee_{n=2}^{\infty}S^n.
\]
Every integral Hurewicz homomorphism $\pi_n(BM)\to H_n(BM;\Z)$,
$n\ge2$, is a split epimorphism. The integral cohomology ring is
$\Z\oplus\bigoplus_{n\ge2}\Z u_n$, with $|u_n|=n$ and
$u_pu_q=0$ for all $p,q\ge2$.
\end{theorem}

\begin{proof}
We give the resolution, cocycles, and actual spherical representatives separately.

\textbf{Step 1: associativity and simple connectivity.}
For $m\ne1$ let $r(m),s(m)\in\{0,1\}$ be its row and column,
respectively, and put $r(y)=s(y)=0$. The product of two nonidentity
elements is $x_{r(m),s(m')}$. Both associations of three such elements
are therefore $x_{r(m),s(m'')}$. Products involving $1$ are immediate.
The four band elements are idempotent and hence become $1$ in the
universal group. The relation $ya=a$ then makes $y$ trivial too.
Thus $U(M)=1$ and $BM$ is simply connected.

\textbf{Step 2: homology in every degree, not just a bar-complex range.}
Put $R=\Z M$ and $P=Ra=\Z\{a,c\}=Rc$. The augmentation of $P$
is denoted $\epsilon$. Define right-multiplication maps
\[
 D:R\longrightarrow P,\quad v\longmapsto v(a-c),
 \qquad
 A:P\oplus R\longrightarrow R,\quad (u,v)\longmapsto ub+vy,
\]
\[
 B:P\oplus R\longrightarrow R,\quad
 (u,v)\longmapsto u(c-d)+v(a-y).
\]
These are homomorphisms of left $R$-modules. All modules displayed are
projective: $a$ is idempotent and $R$ is free. The resolution is
\[
 \cdots\xrightarrow{(0,A)}P\oplus R
 \xrightarrow{(0,B)}P\oplus R
 \xrightarrow{A}R\xrightarrow{D}P
 \xrightarrow{\epsilon}\Z\longrightarrow0,
\]
where the tail alternates $(0,A)$ and $(0,B)$.
Here $\ker\epsilon=\Z(a-c)=\operatorname{im}D$, and
$\ker D=\Z\{a,b,c,d,y\}=\operatorname{im}A$.
The summands $ub$ and $vy$ of $A$ have disjoint supports
$\{b,d\}$ and $\{a,c,y\}$. It follows directly that
\[
 \ker A=0\oplus\Z\{a-b,c-d,a-y\}=0\oplus\operatorname{im}B.
\]
The three displayed generators are independent; the first two come from
$u(c-d)$ and the last from the coefficient of $1$ in $v(a-y)$.
Consequently
\[
 \ker B=0\oplus\Z\{a,b,c,d,y\}=0\oplus\operatorname{im}A.
\]
This proves exactness of the entire infinite resolution. After tensoring
with the trivial right $R$-module $\Z$, $D$ becomes zero, $A$ becomes
$(1\;1)$, and $B$ becomes zero. In higher even degrees the nonzero
map has matrix $\left(\begin{smallmatrix}0&0\\1&1\end{smallmatrix}\right)$.
Thus $H_1(BM;\Z)=0$ and $H_n(BM;\Z)=\Z$ for every $n\ge2$.
This independently confirms the repeating resolution in
\cite[Example 5.2]{sweeney2025}.

\textbf{Step 3: primitive integral cycles.}
Extend $r,s$ by zero at $1$, and put $v(m)=1$ if $m=y$ and $v(m)=0$
otherwise. For $n\ge2$ define the normalized integral cochain
\[
 u_n(m_1,\ldots,m_n)
   =r(m_1)v(m_2)\cdots v(m_{n-1})s(m_n).
\]
For nonidentity $m,m'$ we have
$r(m'm)=r(m')$, $s(m'm)=s(m)$ and $v(m'm)=0$.
In $\delta u_n$, the first two terms cancel, as do the last two;
every remaining term has a factor $v(m'm)=0$. For $n=2$ there
are no remaining terms. Hence $\delta u_n=0$ integrally.

In the multilinearly extended normalized bar group set
\[
 z_n=(c-a)\otimes(y-a)^{\otimes(n-2)}\otimes(b-a).
\]
Every chronological product of adjacent chosen letters is $a$,
independent of either binary choice. Thus the inner faces in
$\partial z_n$ cancel in pairs; the outer faces do too. Hence
$\partial z_n=0$, and $\langle u_n,z_n\rangle=1$.
By Step 2, $[z_n]$ generates $H_n(BM;\Z)$.

\textbf{Step 4: finite spheres which map to these cycles.}
Let $C_n$ have objects $0,\ldots,n$, two arrows $i-1\to i$ for
each adjacent pair, one arrow $i\to j$ for $j-i\ge2$, and no
other nonidentity arrows. All compositions of nonidentity arrows
are the unique arrow across the resulting interval. This defines
a finite category. Write $X_n=BC_n$; $X_0$ is a point and
$X_1$ is a circle. We claim
\[
 X_n\simeq\bigvee^{F_n}S^n,\qquad
 F_0=0,\quad F_1=1,\quad F_n=F_{n-1}+F_{n-2}.
\]
To verify the induction without a homology-only inference, append the
sink $n$ to $C_{n-1}$. The category of arrows from $C_{n-1}$ to $n$
is exactly two copies of $C_{n-1}$ glued along $C_{n-2}$.
Indeed the object $n-1$ occurs twice and every earlier object once;
the arrows in either copy and their compositions are unchanged.
The forgetful functor folds the copies. Decomposing the nerve into
simplices ending at $n$ and those not ending there identifies $X_n$
with the mapping cone of this fold. Since all inclusions are CW
subcomplex inclusions, this gives
\[
 X_n\simeq\operatorname{hocofib}
 (X_{n-1}\cup_{X_{n-2}}X_{n-1}\longrightarrow X_{n-1})
 \simeq\Sigma(X_{n-1}/X_{n-2}).
\]
The second equivalence follows by contracting one of the two cone copies
of $X_{n-1}$ in the mapping-cone description; what remains is the
suspension of the quotient. By induction the inclusion
$X_{n-2}\to X_{n-1}$ is null-homotopic: its domain is a wedge of
$(n-2)$-spheres and its target is $(n-2)$-connected. For $n=2$
the domain is a point. Hence
$X_{n-1}/X_{n-2}\simeq X_{n-1}\vee\Sigma X_{n-2}$,
which proves the claimed recurrence and homotopy equivalence.

Define an actual functor $C_n\to M$ by mapping all objects to the
single object, each long arrow to $a$, the adjacent pairs to
$(a,c)$ at the beginning, $(a,y)$ at each interior step, and
$(a,b)$ at the end. Composable products are $a$, so the functor
axioms hold. The signed sum of the $2^n$ top simplices, with
coefficient $(-1)^{n-\sum\epsilon_i}$ for the binary choices
$\epsilon_i$, is a cycle in $X_n$ and maps to $z_n$.
The preceding wedge identification makes the Hurewicz map
$\pi_n(X_n)\to H_n(X_n;\Z)$ an isomorphism. Thus this cycle
has an actual sphere representative, whose composite $S^n\to BM$
represents the primitive generator $[z_n]$.

\textbf{Step 5: the equivalence and ring.}
Wedge the chosen sphere maps for all $n\ge2$. The resulting map
$\bigvee_{n\ge2}S^n\to BM$ is a homology isomorphism of simply
connected CW complexes and therefore a homotopy equivalence.
There is also a direct cochain verification of the ring assertion.
Put $d(m)=r(m)s(m)$ and, for $p,q\ge2$, set
\[
 w_{p,q}=r\otimes v^{\otimes(p-2)}\otimes d
          \otimes v^{\otimes(q-2)}\otimes s.
\]
The same endpoint cancellations apply to its coboundary.
The only surviving internal face replaces $d(m'm)$ by
$r(m')s(m)$, giving
$\delta w_{p,q}=(-1)^p u_p\smile u_q$.
Thus every positive-degree product vanishes over $\Z$, and the
same explicit primitives work over every coefficient ring.
\end{proof}

\begin{corollary}\label{cor:E-uncountable}
There is a finite monoid, of order six including its identity, for which
$[BM,S^2]$ has cardinality $2^{\aleph_0}$. In particular the assertion
that $[BC,S^2]$ is countable for every finite category is false.
\end{corollary}
\begin{proof}
Theorem~\ref{thm:E-wedge} gives
$[BM,S^2]_*\cong\prod_{n\ge2}\pi_n(S^2)$.
Pointed and unpointed classes agree because the target is simply connected.
Each factor is countable, and infinitely many factors are nonzero
(for example the classical $p$-torsion in $\pi_{2p}(S^3)$,
transported by the Hopf fibration for $2p\ge3$; the nonzero elements
are constructed in \cite[Example 4L.6, p.~495]{hatcher2002}).
Choosing zero or one nonzero element in infinitely many distinct factors
gives the lower bound $2^{\aleph_0}$; a countable product of countable
sets gives the upper bound.
\end{proof}

\subsection{Regular monoids: the exact rational bound}
\begin{proposition}\label{prop:E-regular}
Let $M$ be a finite regular monoid with $j$ Green $\mathcal J$-classes.
Then $H_n(BM;\Q)=0$ for $n>2(j-1)$.
Consequently none of $K(\Z,2)$, $\Omega S^3$, $\Sigma\CP^\infty$
or $\bigvee_{n\ge2}S^n$ is the classifying space of a finite regular monoid.
\end{proposition}
\begin{proof}
Adjoin a new zero to $M$, obtaining $M^0$. Its contracted algebra
$\Q_0[M^0]$ is canonically $\Q M$; the new zero is the only extra
$\mathcal J$-class. Apply \cite[Theorem 5.13]{martiniansteinberg2025}
with $K=\Q$ and a principal series for $M^0$. There are $j$ nonzero
principal factors. Every maximal subgroup is finite, so its order is
a unit in $\Q$, and $\operatorname{gl.dim}\Q=0$. The asserted
bound is therefore $\operatorname{gl.dim}\Q M\le2(j-1)$.
The normalized bar resolution identifies $H_n(BM;\Q)$ with
$\Tor_n^{\Q M}(\Q,\Q)$, which vanishes above global dimension.
Each listed space has nonzero rational homology in infinitely many
degrees, so none can satisfy this bound. This checks all hypotheses
of the cited theorem and avoids the unverified memory citation to Putcha.
\end{proof}

\subsection{The central-cocycle reformulation}
\begin{proposition}\label{prop:E-cocycle}
Let $M$ be a discrete monoid. For a normalized integral monoid two-cocycle
$u:M\times M\to\Z$, let $E_u=M\times\Z$ with
\[
 (a,r)(b,s)=(ab,r+s+u(a,b)).
\]
There is a homotopy-fiber sequence
\[
 B\Z\longrightarrow BE_u\longrightarrow BM
\]
classified by $[u]\in H^2(BM;\Z)$, up to the choice of sign for the fiber
generator. Consequently $BM\simeq K(\Z,2)$ if and only if some such $E_u$
has contractible classifying space. The extra hypothesis $U(M)=1$ in R9(H2)
is necessary but redundant in this reformulation. For finite $M$, neither
condition occurs, by Corollary~\ref{A:rational-nilpotence}.
\end{proposition}
\begin{proof}
Associativity is exactly the cocycle equation
\[
 u(a,b)+u(ab,c)=u(b,c)+u(a,bc),
\]
and normalization makes $(1,0)$ the identity. The projection $E_u\to M$
is a functor between one-object categories. Its categorical fiber is the
group $\{1\}\times\Z$. Every arrow $m$ of the base has the lift $(m,0)$;
this is opcartesian because an arrow over $nm$ factors uniquely through it:
\[
 (nm,t)=(n,t-u(n,m))(m,0).
\]
Transport on the fiber is the identity on $\Z$, since the fiber is central.
All transport functors are therefore equivalences. Quillen's homotopy-fiber
criterion, or the Grothendieck-construction theorem applied to this groupoid
opfibration, gives the displayed homotopy-fiber sequence.

The cleavage $(m,0)$ has composition defect $(1,u(a,b))$.
Changing lifts by an integral function $v:M\to\Z$ changes this defect
by the coboundary of $v$. On a two-simplex $(a,b)$ of $NM$, the obstruction
to extending the chosen fiber trivializations is exactly $u(a,b)$.
Thus the characteristic class of this $K(\Z,1)$-fibration is $[u]$.
This also proves the claimed classification directly from the normalized
nerve cochain complex, with no assumption that the base monoid is a group.

If $BM\simeq K(\Z,2)$, choose a normalized cocycle representing a generator
of $H^2(BM;\Z)$. The boundary homomorphism
$\pi_2(BM)\to\pi_1(B\Z)$ in this fiber sequence is evaluation on this
class and is an isomorphism. Its long exact sequence makes every positive
homotopy group of $BE_u$ zero. The realization is a connected CW complex,
so it is contractible. Conversely, if $BE_u$ is contractible, the same
long exact sequence gives $\pi_2(BM)=\Z$ and all other positive homotopy
groups zero. Hence $BM\simeq K(\Z,2)$ and in particular $U(M)=1$.
This confirms the sketch in R9(H2) and supplies its missing comparison and
classification maps.
\end{proof}

\subsection{Step-by-step check of R10, Proposition 4.1}
\begin{proposition}\label{prop:E-R10-audit}
Let $X$ be a connected CW complex with $H^1(X;\Z)=0$. If $S^n$ is a
homotopy retract of $X$ for infinitely many distinct $n$ with
$\pi_n(S^3)\ne0$, then both $[X,S^3]$ and $[X,S^2]$ are uncountable.
The proof of \emph{OWN\_RESULTS}, Proposition 4.1, is correct.
\end{proposition}
\begin{proof}
\textbf{Step 1: the cellular choices.} Choose maps
$s_n:S^n\to X$ and $f_n:X\to S^n$ with $f_ns_n$ of degree $\pm1$.
The spheres in question have $n\ge3$. By cellular approximation, with
$S^n$ given its zero-cell and $n$-cell structure, arrange that $s_n$ lands
in $X^{(n)}$ and that $f_n$ is constant on $X^{(n-1)}$.
These replacements preserve the degree. Choose a based nonzero class
$a_n:S^n\to S^3$, where $S^3$ is the topological group of unit quaternions,
and set $g_n=a_nf_n$.

\textbf{Step 2: the infinite product really is a map.}
For each subset $A$ of the infinite index set define
$g_A(x)=\prod_{n\in A}g_n(x)$ in increasing order.
On $X^{(k)}$ only the finitely many indices $n\le k$ can contribute.
Thus its restriction to every closed cell is a finite product of continuous
maps. The weak topology of a CW complex proves continuity on all of $X$;
no local finiteness of the cell structure is needed.

\textbf{Step 3: separation of two subsets.}
For $A\ne A'$ let $n$ be their first differing index and suppose $n\in A$.
On $X^{(n)}$, the two products are $hg_n$ and $h$, where $h$ is the common
product over smaller indices. If the maps on $X$ were homotopic, restriction
and pointwise multiplication by $h^{-1}$ would make
$g_n|_{X^{(n)}}$ null-homotopic. Composing with $s_n$ would kill
$a_nf_ns_n$, which represents $\pm[a_n]\ne0$ in $\pi_n(S^3)$.
The target is simply connected, so free versus based homotopy introduces
no ambiguity. This is a contradiction. Thus the power set of a countably
infinite subset injects into $[X,S^3]$.

\textbf{Step 4: transfer to the two-sphere.}
For the principal Hopf bundle $S^1\to S^3\to S^2$, two maps from $X$ to
$S^3$ whose projected maps are homotopic become two lifts of the same map
after lifting that homotopy. Two such lifts differ by a continuous map
$X\to S^1$. Its homotopy class belongs to $H^1(X;\Z)$, which is zero.
Multiplication by a null-homotopic map gives homotopic lifts. Therefore
$[X,S^3]\to[X,S^2]$ is injective. This proves precisely the last step of
the supplied argument, and completes its audit.
\end{proof}

\subsection{A certified small-order obstruction for $\Sigma\CP^\infty$}
\label{E:sigma-small-section}

\begin{proposition}[Published homology prefilter and independently checked operations]
\label{E:sigma-small}
Using the complete corner-irreducible monoid census and its published
integral homology through degree seven in orders at most nine and
through degree five in order ten, no monoid of order at most ten
has classifying space homotopy equivalent to $\Sigma\CP^\infty$.
The first monoid order not excluded by this calculation is eleven.
The general exclusion in Corollary~\ref{E:Sigma-CP-infty-excluded}
below is independent of this published prefilter and covers arbitrary
finite monoids and categories.
\end{proposition}

\begin{proof}
Corner reduction preserves the classifying space and does not increase
the order, as in Proposition~\ref{prop:B-Moore-minimum}. Passing to the
opposite monoid also preserves the classifying space: reverse each
simplex of its nerve. Consequently it suffices to examine the complete
corner-irreducible representatives, up to isomorphism and opposition,
in the source census \cite{SweeneyEnumeration2026B}.

\textbf{The source-dependent prefilter.}
The published file \path{fetched/Sweeney_maxorder9_maxdim7.hdf5}
contains the integral homology profiles through degree seven for all
3960 representatives of orders one through nine. Its index counts by
order are
\[
 1,1,1,2,2,7,25,256,3665.
\]
There are no entries marked uncomputed. A necessary profile for
$\Sigma\CP^\infty$ is
\[
 (H_0,H_1,\ldots,H_7)
   =(\Z,0,0,\Z,0,\Z,0,\Z).
\]
Exactly seven order-eight tables and 66 order-nine tables have this
published profile; none of the smaller tables does. The zero-based
order-eight indices are $24,51,59,69,106,135,148$.
The 66 order-nine indices and all 73 multiplication tables are saved
in \path{code/outputs/a_sigma_CP_infty_candidates9.json}.
The reproducible extraction check
\path{code/a_sigma_CP_prefilter.py} checks every source index, the
absence of unknown profiles, exact equality of the extracted tables
with the subsequent operation certificates, and SHA-256 hashes of
the three source files used. Completeness of the census and correctness
of the published integral homology used to discard the other 3887
tables are external source dependencies; they are not reclassified
as independent computations here.

The separate published file
\path{fetched/Sweeney_maxorder10_maxdim5.hdf5} provides integral
homology through degree five for all 71\,916 order-ten core
representatives, again with no uncomputed entries. The necessary
profile $(\Z,0,0,\Z,0,\Z)$ retains exactly 2407 tables. The program
\path{code/a_sigma_CP_steenrod10.py} reproduces this full prefilter,
records its complete selected-index list and both source-file hashes
in \path{code/outputs/a_sigma_CP_prefilter10.json}, and checks all
retained tables. The homology source was also matched to the exact
Git blob recorded in the fetched upstream repository tree; this
provenance check is saved in
\path{code/outputs/a_sigma_CP_source10.json}. These published order-ten
homology values and census completeness are additional explicit
source dependencies.

\textbf{The independently computed obstruction.}
For every retained table the program
\path{code/a_sigma_CP_steenrod.py} first checks the identity and all
associativity equations, then uses the normalized bar differential
over $\mathbb F_2$ to compute $H^3$ independently. It has dimension
one in every case. The same checks in the order-ten runner give
dimension one for all 2407 additional tables. All 73 smaller monoids
and 2316 of the order-ten monoids are independently checked to be
aperiodic, although aperiodicity is not required for the obstruction.
Choose a normalized three-cocycle $x$ representing the unique
nonzero class.

For this computation we use algebraic bar order, so the internal
face replaces consecutive entries $a,b$ by $ab$. This is the reverse
of chronological nerve order used in Theorem~\ref{thm:E-wedge};
reversal identifies the two conventions. Extend $x$ by zero whenever
one argument is the identity. The normalized cup-one formula gives
\[
\begin{split}
 (x\smile_1x)(a,b,c,d,e)
   ={}&x(abc,d,e)x(a,b,c)\\
      &+x(a,bcd,e)x(b,c,d)\\
      &+x(a,b,cde)x(c,d,e).
\end{split}
\]
This cocycle represents $\mathrm{Sq}^2[x]$; the displayed insertion
formula is also the bar formula in \cite[Theorem~3.7 and its preceding
definition]{lodder2018}. Exact binary linear algebra produces a
normalized four-cochain $v$ satisfying
\[
                 \delta v=x\smile_1x
\]
for every one of the 2480 retained tables. The output retains $x$, $x\smile_1x$, and
$v$ as complete bit vectors in the lexicographic normalized-word
bases. The program verifies the equality on every normalized
five-tuple: $7^5$ equations for each order-eight table and $8^5$
for each order-nine table, for a total of 2\,280\,337 equations in
those orders. The order-ten run checks another
$2407\cdot9^5=142\,130\,943$ equations. Its complete per-table
certificates are saved in
\path{code/outputs/a_sigma_CP_steenrod10.jsonl}; the completed summary
records all 2407 distinct source indices and no nonzero operation.
The combined computation checks 144\,411\,280 primitive identities.
Thus $\mathrm{Sq}^2:H^3(BM;\mathbb F_2)\to H^5(BM;\mathbb F_2)$
is zero in all 2480 cases. As an additional control the same code
recovers $\mathrm{Sq}^2(w^3)=w^5\ne0$ for $BC_2$.

Finally, if $c\in H^2(\CP^\infty;\mathbb F_2)$ is the standard
generator and $\sigma$ denotes cohomological suspension, stability
and the top-square identity give
\[
 \mathrm{Sq}^2(\sigma c)
       =\sigma\mathrm{Sq}^2(c)=\sigma(c^2)\ne0.
\]
This contradicts the independently certified vanishing for every
candidate and proves the stated exclusion with precisely the
source dependencies recorded above.
\end{proof}

\begin{remark}\label{E:sigma-open}
Cup-product nilpotence does not obstruct a suspension, whose
positive-degree products vanish.  The finite operation certificates
above do not by themselves establish an all-order assertion.
Corollary~\ref{E:general-good-prime-LF} supplies the separate general
Steenrod-module theorem, and
Corollary~\ref{E:Sigma-CP-infty-excluded} settles the realization
question negatively for both finite monoids and finite categories.
\end{remark}

\subsection{A nonregular family with locally finite Steenrod operations}

\begin{theorem}\label{E:nilpotent-inflations}
Let $N$ be a finite nilpotent semigroup with zero $0$, and let $I,J$
be finite nonempty sets with chosen elements $i_0,j_0$.  Define
\[
 M=\{1\}\amalg(I\times J)\amalg(N\setminus\{0\})
\]
as follows.  The element $1$ is the identity, the set $I\times J$ is a
rectangular-band ideal, and, writing $e=(i_0,j_0)$,
\[
 (i,j)(i',j')=(i,j'),\qquad
 n(i,j)=(i_0,j),\qquad (i,j)n=(i,j_0).
\]
For $n,n'\in N\setminus\{0\}$, their product is their product in $N$
if that is nonzero, and is $e$ otherwise.
Then $M$ is an aperiodic finite monoid, and for every prime $p$ its
cohomology $H^*(BM;\mathbb F_p)$ is locally finite as a module over the
Steenrod algebra.  In particular $BM$ is not homotopy equivalent to
$\Sigma\CP^\infty$.

More precisely, let $Q=N^1$ be the pointed monoid with zero as basepoint
and a freshly adjoined identity, and let
$\epsilon:Q\to S^0=\{*,1\}$ send all nonunits to the basepoint.
Let $X_N$ be the based two-sided simplicial bar
$B(S^0,Q,S^0)$ for this augmentation, with smash products in its levels.
There is a stable equivalence
\[
 \widetilde{\Sigma^\infty BM}\ \simeq
 \bigvee^{(|I|-1)(|J|-1)}\Sigma^2\Sigma^\infty X_N.       \tag{E.20}
\]
Every reduced cohomology class of $X_N$ is pulled back from a finite
CW quotient of $X_N$.
\end{theorem}

\begin{proof}
\emph{The multiplication and the spherical ideal.}
Associativity of three outside factors follows from associativity in
$N$, with its zero replaced by $e$.  The element $e$ absorbs every
outside factor.  With one or two band factors, the displayed row and
column rules give the same outer coordinates in either association.
These checks cover all triples.  An outside element has a power equal
to $e$, and every band element is idempotent, so all eventual cyclic
groups are trivial.  Thus $M$ is aperiodic.

We give the stable argument explicitly because an assertion only about
the ordinary integral monoid algebra would not control Steenrod
operations.  Write $\mathbb S$ for the sphere spectrum and put
\[
 R=\Sigma^\infty M_+,\qquad
 K=\Sigma^\infty(I\times J)_+,\qquad
 A=\Sigma^\infty Q.
\]
The multiplication of a discrete monoid makes its suspension spectrum
an associative ring spectrum.  The ideal inclusion and quotient give
a cofiber sequence of $R$-bimodules
\[
 K\longrightarrow R\longrightarrow A.                 \tag{E.21}
\]
Here the quotient identifies the entire band with the basepoint; it
does not identify the outside elements with one another.

The left module $Re$ and right module $eR$ are literal retracts of
the regular modules, using multiplication by the idempotent $e$.
Their underlying spectra are the wedges on $I$ and $J$, respectively,
and $eRe=\mathbb S$.  Multiplication induces an equivalence
\[
 Re\mathbin{\wedge}_{\mathbb S}eR\ \xrightarrow{\ \cong\ }K
                                                        \tag{E.22}
\]
of bimodules: on the indexing sets it is the bijection
$((i,j_0),(i_0,j))\mapsto(i,j)$.
Derived tensor products with these representable retracts are their
ordinary idempotent summands.  In particular $Ae=eA=0$.  Tensoring
(E.21) with $A$ and using (E.22) proves
\[
 A\mathbin{\wedge}^{\mathbb L}_R A\simeq A.             \tag{E.23}
\]
Consequently restriction from $A$-modules to $R$-modules is fully
faithful: its extension--restriction counit is an equivalence by
(E.23) and associativity of derived tensor products.

\emph{The two augmentation fibers.}
Regard $\mathbb S$ as a trivial left or right $R$-module, so every
element of $M$ acts as the identity.  Let
\[
 V=\operatorname{fib}(Re\longrightarrow\mathbb S),\qquad
 W=\operatorname{fib}(eR\longrightarrow\mathbb S),        \tag{E.24}
\]
where the maps fold the indicated wedges.  Applying $e$ to the first
fiber gives the fiber of the literal identity
$\mathbb S\to\mathbb S$, so $eV\simeq0$; similarly $We\simeq0$.
This is a vanishing of spectra with their module structure, not just
of their degree-zero homotopy groups.  Equation (E.22) therefore gives
$K\wedge_R^{\mathbb L}V\simeq0$.  Tensoring (E.21) with $V$ gives
\[
 V\simeq A\mathbin{\wedge}^{\mathbb L}_R V.
\]
It supplies a coherent $A$-module structure on $V$.  The analogous
argument gives one on $W$.

The usual simplicial bar is a model for derived tensor products, and
in particular
$\mathbb S\wedge_R^{\mathbb L}\mathbb S\simeq\Sigma^\infty_+BM$.
Tensor the first fiber sequence (E.24) with the trivial right module.
Its middle term is $\mathbb S e=\mathbb S$, and the resulting map to
$\Sigma^\infty_+BM$ is split by the augmentation to $\mathbb S$.
Taking the reduced summand gives
\[
 \widetilde{\Sigma^\infty BM}
       \simeq\Sigma(\mathbb S\wedge_R^{\mathbb L}V).
\]
Now tensor the second sequence (E.24) with $V$.  Its middle term is
$eV\simeq0$, whence
\[
 \widetilde{\Sigma^\infty BM}
       \simeq\Sigma^2(W\wedge_R^{\mathbb L}V)
       \simeq\Sigma^2(W\wedge_A^{\mathbb L}V).           \tag{E.25}
\]
The last equivalence follows from (E.23) since both modules descend
to $A$.

\emph{Identification of the modules, including coherence.}
Let $T=\{1,e\}\amalg(N\setminus\{0\})\subset M$.
The element $e$ is a zero in $T$.  Put $P=\Sigma^\infty T_+$.
The same one-element ideal calculation gives
$A\wedge_P^{\mathbb L}A\simeq A$, so restriction of $A$-modules to
$P$-modules is fully faithful as well.

On restriction to $T$, the row $(i_0,j_0)$ defines a module section
$\mathbb S\to Re$ of its augmentation: it is fixed by every element
of $T$.  In a stable category the fiber of a split epimorphism is the
cofiber of its section.  Here that cofiber is the strict pointed
permutation quotient $\Sigma^\infty(I/\{i_0\})$.
Every nonunit of $T$ sends all rows to $i_0$, hence acts as the zero
map on this quotient.  Thus, as a $P$-module, it is precisely a wedge
of $|I|-1$ copies of the sphere module defined by $\epsilon:A\to\mathbb S$.
Both it and $V$ are restrictions of $A$-modules.  Full faithfulness
therefore lifts their equivalence to an equivalence of $A$-modules.
The column argument gives, with the proper right action,
\[
 V\simeq\bigvee^{|I|-1}\mathbb S_\epsilon,
 \qquad W\simeq\bigvee^{|J|-1}\mathbb S_\epsilon.
\]
This argument uses actual module maps and the two quotient equivalences;
no unproved lift of a degree-zero idempotent or matrix is used.
The based simplicial bar identifies
$\mathbb S_\epsilon\wedge_A^{\mathbb L}\mathbb S_\epsilon$
with $\Sigma^\infty X_N$.  Substitution in (E.25) proves (E.20).

\emph{Finite quotients detecting every cohomology class.}
Choose $\ell\geq2$ with $N^\ell=\{0\}$.  For $a\ne0$ set
\[
 w(a)=\max\{r:a\in N^r\}.
\]
Then $1\leq w(a)\leq\ell-1$, and if $ab\ne0$ then
$w(ab)\geq w(a)+w(b)$.  A nonbasepoint nondegenerate positive-dimensional
simplex of $X_N$ is a word $(a_1,\ldots,a_n)$ of nonzero elements of $N$.
Its two exterior faces are the basepoint.  An interior face multiplies
two adjacent letters, and is the basepoint if their product is zero.
Otherwise it does not decrease the total weight
$\sum_iw(a_i)$.  Degeneracies insert the identity, of weight zero.

For an integer $q\geq1$, the basepoint and all words of total weight
at least $q$ consequently form a based simplicial subspace $F_qX_N$.
Every nondegenerate cell of this subspace has dimension at least
$\lceil q/(\ell-1)\rceil$.  In contrast, every nondegenerate positive
cell of the quotient $X_N/F_qX_N$ has dimension at most $q-1$,
because every letter has weight at least one.  The quotient is a finite
CW complex, since $N$ is finite.  Its nonbasepoint zero-simplex also
causes no problem and is retained in the quotient.

Given a fixed degree $d$, choose $q>(d+1)(\ell-1)$.
The subspace $F_qX_N$ has no positive cells in dimensions at most
$d+1$.  The cofiber sequence shows that
\[
 \widetilde H^d(X_N/F_qX_N;\mathbb F_p)
     \longrightarrow\widetilde H^d(X_N;\mathbb F_p)
\]
is an isomorphism.  Thus every class lifts to a finite quotient.
Naturality of Steenrod operations makes the entire Steenrod orbit
of that class an image of the finite-dimensional cohomology of this
one quotient.  It is therefore finite-dimensional.  The same holds
for (E.20), since Steenrod operations commute with suspension and
stable equivalences.

Finally, if $u\in H^2(\CP^\infty;\mathbb F_p)$ is the generator,
its suspension has nonzero successive operations
$P^{p^{r-1}}\cdots P^1(\sigma u)=\sigma(u^{p^r})$
at every odd prime.  At $p=2$ the explicit sequence is
$Sq^{2^r}\cdots Sq^4 Sq^2(\sigma u)=\sigma(u^{2^r})$ for $r\geq1$,
where the first operation is $Sq^2$ and the intermediate string is
empty when $r=1$.  Its Steenrod orbit is infinite-dimensional,
which proves the asserted exclusion.
\end{proof}

\begin{remark}\label{E:nilpotent-inflations-scope}
The theorem includes the inflated rectangular-band infinite-wedge
example (one nonzero nilpotent letter with square zero).  Its proof
does not assume bounded homology or finite-dimensional homotopy type.
It proves local finiteness for the stated family, not for all aperiodic
monoids or all finite categories.  The next subsection removes the
constant-action hypothesis and allows a finite unit group at suitable
primes.  Further nonzero idempotents in a general Rees quotient
require a separate argument; the controlled-image localization and
multibar induction of Theorem~\ref{E:general-controlled-multibar}
provide it below.
\end{remark}

\subsection{Nilpotent Rees quotients with arbitrary row and column actions}

\begin{proposition}\label{E:band-free-group-bar}
Let $M$ be a finite monoid whose minimal ideal is a rectangular band
$K=I\times J$.  Write $Q=M/K$ for the pointed Rees quotient, with the
entire ideal identified to its zero.  For a finite nonempty set $I$
with chosen element $0$, let $F_I$ be the free group on symbols $x_i$
for $i\ne0$, and set $x_0=1$.  A set map $\sigma:I\to I$ induces the
group homomorphism
\[
 \Phi_\sigma(x_i)=x_{\sigma(0)}^{-1}x_{\sigma(i)}.       \tag{E.26}
\]
The row actions of $M$ on $I$ give a pointed left $Q$-action on $BF_I$,
and the column actions give a pointed right $Q$-action on $BF_J$.
There is a stable equivalence
\[
 \widetilde{\Sigma^\infty BM}\simeq
 \Sigma^\infty B(BF_J,Q,BF_I).                         \tag{E.27}
\]
Here the bar is the \emph{pointed, smash-product} two-sided bar: its
degree-$n$ space is $BF_J\wedge Q^{\wedge n}\wedge BF_I$.
\end{proposition}

\begin{proof}
Use ordinary monoid multiplication, so $m(i,j)=(\sigma_m(i),j)$
and $(i,j)m=(i,\tau_m(j))$.  Thus
$\sigma_{mn}=\sigma_m\sigma_n$ and
$\tau_{mn}=\tau_n\tau_m$.  Substitution in (E.26) gives the literal
identity $\Phi_\sigma\Phi_\rho=\Phi_{\sigma\rho}$:
\[
 \Phi_\sigma(x_{\rho(0)}^{-1}x_{\rho(i)})
 =x_{\sigma\rho(0)}^{-1}x_{\sigma\rho(i)}.
\]
A constant set map induces the trivial group homomorphism.  Every
element of the rectangular band acts constantly on each of $I,J$,
so these pointed actions factor strictly through $Q$.

We give the coherence in the stable identification explicitly.
Put $R=\Sigma^\infty M_+$, $A=\Sigma^\infty Q$, and choose
$e=(0,0)\in K$.  The ideal spectrum is literally
\[
 \Sigma^\infty K_+\cong Re\wedge_{\mathbb S}eR
\]
as an $R$-bimodule, by the bijection
$((i,0),(0,j))\mapsto(i,j)$.  Since $Re,eR$ are idempotent retracts
of regular modules and $eRe=\mathbb S$, the ideal--quotient cofiber
sequence gives
$A\wedge_R^{\mathbb L}A\simeq A$.  Restriction of $A$-modules to
$R$-modules is therefore fully faithful.

Let $V=\operatorname{fib}(Re\to\mathbb S)$ and
$W=\operatorname{fib}(eR\to\mathbb S)$ for the trivial augmentation
modules.  Applying $e$ gives the fiber of the identity of $\mathbb S$,
so $eV=We\simeq0$ as spectra.  Tensoring the ideal--quotient sequence
then supplies coherent equivalences
$V\simeq A\wedge_R^{\mathbb L}V$ and
$W\simeq W\wedge_R^{\mathbb L}A$.  The two augmentation fiber
sequences, exactly as in (E.25), give
\[
 \widetilde{\Sigma^\infty BM}
 \simeq\Sigma^2(W\wedge_A^{\mathbb L}V).               \tag{E.28}
\]

Let $\Gamma_I$ be the graph with two vertices $v_0,v_1$ and one
oriented edge $a_i:v_0\to v_1$ for every $i\in I$, based at $v_1$.
The cofiber of the fold $\Sigma^\infty I_+\to\mathbb S$ is naturally
$\Sigma^\infty\Gamma_I$.  Consequently
$\Sigma V\simeq\Sigma^\infty\Gamma_I$, naturally for all set maps
of $I$.  The same statement holds for columns.

For precision one can implement the next natural equivalence by a
literal equivariant zigzag.  Let $\mathcal G_I$ be the free groupoid
on the directed graph $\Gamma_I$.  There are natural based maps
\[
 \Gamma_I\longrightarrow B\mathcal G_I
       \longleftarrow B\operatorname{Aut}_{\mathcal G_I}(v_1).
                                                               \tag{E.29}
\]
The first map sends an edge to its generating arrow.  The second is
the full subgroupoid inclusion at $v_1$.  Both are homotopy
equivalences: the graph is aspherical and has this fundamental
groupoid, and $\mathcal G_I$ is connected.  Both maps commute with
every set map $I\to I$, because such maps fix the two vertices.
Identifying $x_i=a_0^{-1}a_i$ identifies the last group with $F_I$
and its induced maps with (E.26).  Thus (E.29), not merely a
degree-zero module calculation, gives an $R$-module equivalence
\[
 \Sigma V\simeq\Sigma^\infty BF_I.
\]
The right side has the strict $A$-module structure already described;
full faithfulness of restriction upgrades this to an equivalence of
$A$-modules.  Similarly $\Sigma W\simeq\Sigma^\infty BF_J$ as right
$A$-modules.  Substitute in (E.28).  Suspension spectra preserve
smash products and geometric realizations, so the derived tensor
product of these two suspension spectra is the suspension spectrum
of the pointed two-sided bar.  This proves (E.27).
\end{proof}

\begin{lemma}\label{E:free-subset-intersections}
For $S\subseteq I$, put
$H_S=\langle x_a^{-1}x_b:a,b\in S\rangle\leq F_I$,
with $H_\varnothing=1$.  Then
\[
 \operatorname{im}\Phi_\sigma=H_{\sigma(I)},\qquad
 H_S\cap H_T=H_{S\cap T}.                              \tag{E.30}
\]
Every finite union of the nerve subcomplexes $BH_S\subset BF_I$
has finite CW homotopy type.  The same holds for finite unions of
products or smash products of such subcomplexes, and for cofibers of
inclusions between these unions.
\end{lemma}

\begin{proof}
The first identity follows directly from (E.26): its generators are
$x_{\sigma(0)}^{-1}x_b$ for $b\in\sigma(I)$, and their pairwise
quotients generate all $x_a^{-1}x_b$.  Regard $H_S$ as the image of
the fundamental group, based at $v_1$, of the subgraph with the
edges indexed by $S$.  A reduced edge loop in a graph is unique.
If the same element has reduced representatives using only $S$ and
only $T$, its reduced representative uses only $S\cap T$.
This proves the intersection identity, including the empty case.
For nonempty $S$, the group $H_S$ is free of rank $|S|-1$.

Nerve intersections are literal subgroup intersections, hence any
finite intersection of the displayed subcomplexes is the classifying
space of a finitely generated free group, and has finite graph
homotopy type.  To pass from intersections to unions, induct on the
number of subcomplexes.  The union of two CW subcomplexes is their
homotopy pushout over their intersection; the latter is itself a
union of fewer such intersections.  Homotopy pushouts of spaces of
finite CW homotopy type have finite CW homotopy type, by cellular
approximation of the two maps between finite replacements.
This proves the assertion for unions.  Finite products and smash
products have the same property, and their intersections have the
corresponding product form.  The same induction applies to them.
Finally, the indicated inclusions are CW cofibrations, so their
quotients are homotopy cofibers of maps between spaces of finite
CW homotopy type.
\end{proof}

\begin{theorem}\label{E:nilpotent-Rees-quotient}
Let $M$ be a finite monoid with rectangular-band minimal ideal $K$.
Suppose its pointed Rees quotient is $Q=N^1$, where $N$ is a finite
nilpotent semigroup with zero and the identity is freshly adjoined.
No fixed-row or fixed-column hypothesis is imposed on the outside
actions.  Then, for every prime $p$, $H^*(BM;\mathbb F_p)$ is locally
finite over the Steenrod algebra.  In particular $BM\not\simeq
\Sigma\CP^\infty$.
\end{theorem}



\begin{proof}
Write $X=BF_I$, $Y=BF_J$ with their pointed actions from
Proposition~\ref{E:band-free-group-bar}, and put
$Z=B(Y,Q,X)$.  Choose $\ell\geq2$ with $N^\ell=\{0\}$ and set
\[
 w(a)=\max\{r:a\in N^r\},\qquad a\ne0.
\]
Thus $1\leq w(a)\leq h:=\ell-1$, and nonzero products satisfy
$w(ab)\geq w(a)+w(b)$.  Define descending based subcomplexes
\[
 X^0=X,\qquad
 X^s=\bigcup_{\substack{a\ne0\\w(a)\geq s}}
          B\operatorname{im}\Phi_a\quad(s\geq1),       \tag{E.31}
\]
where an empty union means the basepoint.  Define $Y^s$ in the same
way for the right action.  In particular $X^{h+1}=Y^{h+1}=*$.
For a nonzero nonunit $b$, the action maps satisfy
\[
 bX^s\subseteq X^{s+w(b)},\qquad
 Y^s b\subseteq Y^{s+w(b)}.                             \tag{E.32}
\]
For $s=0$ this is the definition of the image of $b$.  For $s>0$ it
follows by composing $b$ with the actor defining an image in (E.31)
and using superadditivity of $w$; a zero product gives the basepoint.
Every subcomplex in (E.31), and every finite union or intersection
needed below, has finite CW homotopy type by
Lemma~\ref{E:free-subset-intersections}.

A nondegenerate middle word in bar degree $n$ is
$\mathbf a=(a_1,\ldots,a_n)$ with all $a_i\ne0,1$; write
$w(\mathbf a)=\sum_iw(a_i)$.  In the summand $Y\wedge X$ indexed by
this word, define the level-$L$ subspace to be
\[
 \bigcup_{s+t+w(\mathbf a)\geq L}Y^s\wedge X^t,
 \qquad 0\leq s,t\leq h.                              \tag{E.33}
\]
Give the identity letter weight zero, so this definition extends to
degenerate words.  The basepoint and all these subspaces define a
based simplicial subspace $F_LZ_\bullet$ of the pointed bar.
Indeed an inner face either gives the basepoint or replaces two
letters by one of at least their total weight.  An end face applies
the deleted letter to an endpoint and, by (E.32), transfers its
weight to the endpoint filtration.  Degeneracies insert a zero-weight
identity.  Thus every simplicial operator preserves (E.33).

There are two complementary finiteness estimates.  First, (E.33)
is the basepoint whenever $(n+2)h<L$.  In a CW realization, a
nonbasepoint bar simplex in degree $n$ contributes cells only in
dimensions at least $n$.  Hence $F_LZ:=|F_LZ_\bullet|$ has no
nonbasepoint cells below degree $\lceil L/h\rceil-2$.
In particular its connectivity tends to infinity with $L$.

Second, if a nondegenerate middle word has length $n\geq L$, its
weight is at least $L$ and (E.33) contains its entire summand, using
$s=t=0$.  The quotient simplicial space therefore has no
nondegenerate middle words in degrees $n\geq L$.  In every smaller
degree there are only finitely many words.  Each remaining summand
is the cofiber of (E.33) in $Y\wedge X$.  Formula (E.33) is a finite
union of smash products of the finite unions in (E.31), and its
intersections have the same form.  By the lemma, this cofiber has
finite CW homotopy type.  The geometric realization is built in
finitely many stages from finitely many such spaces, using their
latching cofibrations and products with finite simplices.  Thus
\[
 Z/F_LZ\quad\hbox{has finite CW homotopy type}.          \tag{E.34}
\]
This is a homotopy-type assertion: the nerves $BF_I$ themselves
usually have infinitely many cells.

Fix a cohomological degree $d$.  Choosing $L$ so that
$\lceil L/h\rceil-2>d+1$ makes the quotient map induce an
isomorphism on reduced degree-$d$ cohomology.  Hence any class in
$\widetilde H^d(Z;\mathbb F_p)$ is pulled back from the cohomology
of one space (E.34), whose total cohomology is finite-dimensional.
Naturality makes every iterated Steenrod operation on that class
the pullback of the corresponding operation on this one lift.
Its entire Steenrod orbit is therefore finite-dimensional.
Proposition~\ref{E:band-free-group-bar} transfers this conclusion
to $BM$.  The explicit nonzero operation sequences on the suspended
generator of $H^2(\CP^\infty;\mathbb F_p)$, recorded in
Theorem~\ref{E:nilpotent-inflations}, exclude $\Sigma\CP^\infty$.
\end{proof}

\begin{remark}\label{E:nilpotent-Rees-scope}
This theorem extends the constant-action inflation family.  Its proof
alone does not give Steenrod local finiteness for all finite aperiodic monoids.
Nonzero idempotents in the Rees quotient prevent the positive weight
used above.  Extending the argument to further regular ideal layers
requires additional coherent control of the endpoint module spectra;
an ordinary monoid-algebra composition series does not provide it.
The strictification in
Lemma~\ref{E:regular-controlled-localization} supplies this control,
and Theorem~\ref{E:general-controlled-multibar} proves the general
result below.
\end{remark}

\begin{lemma}\label{E:relative-bar-good-prime}
Let $Q$ be a finite pointed monoid with group of units $G$, and let
$Y,X$ be pointed simplicial sets with right and left $Q$-actions.
Suppose the action of $G$ is by simplicial automorphisms.  Form the
ordinary pointed relative bar
\[
 C_n=Y\wedge_G Q\wedge_G\cdots\wedge_G Q\wedge_G X
                                                               \tag{E.35}
\]
with $n$ middle copies of $Q$.  The balancing is ordinary orbit
balancing, not derived balancing.  If $p\nmid|G|$, the canonical map
\[
 B(Y,Q,X)\longrightarrow |C_\bullet|                  \tag{E.36}
\]
is an $\mathbb F_p$-homology equivalence.  It is an actual map of
pointed spaces, and hence its cohomology isomorphism commutes with
all Steenrod operations.
\end{lemma}

\begin{proof}
Put $k=\mathbb F_p$.  Degreewise orbit balancing is compatible with
the bar faces and degeneracies, and defines (E.36).  To verify its
homology, retain the simplicial direction inside $X,Y$ and the
separate outer bar direction.  For a simplicial set $T$ with a finite
group $H$ acting simplicially, normalized chains satisfy
\[
 N\widetilde k[T/H]\cong (N\widetilde k[T])_H.          \tag{E.37}
\]
Indeed the nondegenerate simplices of $T/H$ are precisely the orbits
of nondegenerate simplices of $T$: if an orbit is degenerate, any
representative is a translate of a degenerate simplex and is itself
degenerate.  This also holds for pointed quotients.

The $n$-fold balancing in (E.35) is the orbit quotient by $G^{n+1}$,
acting by
\[
 (y,a_1,\ldots,a_n,x)\longmapsto
 (yg_0,g_0^{-1}a_1g_1,\ldots,g_{n-1}^{-1}a_ng_n,g_n^{-1}x).
\]
The natural Eilenberg--Zilber comparison for products and smash
products is equivariant.  Since $kG$ is semisimple, coinvariants are
exact and preserve this chain equivalence.  Thus the inner chains
of (E.35) are naturally modeled by
\[
 N\widetilde k[Y]\otimes_{kG} k[Q]\otimes_{kG}\cdots
       \otimes_{kG} k[Q]\otimes_{kG}N\widetilde k[X],   \tag{E.38}
\]
where $k[Q]$ is the contracted algebra, so the zero of $Q$ is zero.
The absolute bar has the same expression with tensor products over
$k$ in place of $kG$.

For completeness, the relative bar computes the same derived tensor
product as the absolute bar.  Its augmented resolution of a left
$k[Q]$-module $L$ has terms
$k[Q]\otimes_{kG}\cdots\otimes_{kG} k[Q]\otimes_{kG}L$.
The usual extra degeneracy inserting the unit contracts the augmented
complex after forgetting the left action.  Every term is projective
as a left $k[Q]$-module, since its last tensor factor is a $kG$-module
and all $kG$-modules are projective.  The same argument degree by
degree applies to bounded-below complexes of modules, as here;
the resulting double complexes are in the first quadrant.
Consequently the quotient comparison from the absolute to the
relative bar is a quasi-isomorphism.  Applied to (E.38) this proves
(E.36), including the stated naturality.
\end{proof}

\begin{theorem}\label{E:nilpotent-Rees-with-units}
Let $M$ have rectangular-band minimal ideal $K$, and suppose
$Q=M/K$ is the disjoint union of a finite group of units $G$ and a
nilpotent ideal $N$ with zero.  At every prime $p\nmid|G|$,
$H^*(BM;\mathbb F_p)$ is locally finite over the Steenrod algebra.
In particular $BM\not\simeq\Sigma\CP^\infty$.
\end{theorem}

\begin{proof}
Apply Proposition~\ref{E:band-free-group-bar} and then
Lemma~\ref{E:relative-bar-good-prime}.  It suffices to prove local
finiteness for the relative bar (E.35), with $Y=BF_J,X=BF_I$.
For $a\in N\setminus\{0\}$ use the weight
$w(a)=\max\{r:a\in N^r\}$ and give units weight zero.
Multiplication on either side by a unit preserves $N^r$, and hence
preserves $w$.  It also permutes the collection of image subgroups
in (E.31), so the endpoint filtrations $X^s,Y^s$ are $G$-invariant.
The filtration (E.33) therefore descends to the balanced quotient.
End faces still trade letter weight for endpoint filtration, and
inner faces do not decrease weight.

In the relative bar a middle unit factor is literally degenerate:
it can be absorbed into an adjacent factor by balancing and then
is the inserted identity.  The nondegenerate middle factors are
therefore exactly letters of $N\setminus\{0\}$, modulo the finite
balancing action of $G^{n+1}$.  The high-filtration connectivity
estimate in Theorem~\ref{E:nilpotent-Rees-quotient} applies unchanged,
and the quotient has no nondegenerate middle levels $n\geq L$.

Before taking the finite balancing quotient, every level quotient
has finite CW homotopy type by Lemma~\ref{E:free-subset-intersections}.
Equation (E.37) and exactness of $k[G^{n+1}]$-coinvariants show that
its orbit quotient has finite-dimensional total $k$-homology.
The quotient relative bar has only finitely many nondegenerate
middle levels.  Its skeletal homology spectral sequence consequently
shows that its total $k$-homology, and hence total $k$-cohomology, is
finite-dimensional.  No assertion of integral finite CW type for
these orbit quotients is needed.

As before, each fixed-degree cohomology class lifts to one such
high-filtration quotient.  Its Steenrod orbit is contained in the
image of the finite-dimensional cohomology of that quotient.
This proves local finiteness, and the suspended polynomial generator
excludes $\Sigma\CP^\infty$.
\end{proof}

\begin{corollary}\label{E:two-regular-classes}
Let $M$ be a finite monoid whose only regular $\mathcal J$-classes
are its minimal ideal and its group of units (allowing these to
coincide).  Let $H$ be a maximal subgroup in the minimal ideal and
$G$ the unit group.  If $p\nmid |G||H|$, then
$H^*(BM;\mathbb F_p)$ is locally finite over the Steenrod algebra.
Consequently no such monoid has classifying space
$\Sigma\CP^\infty$.
\end{corollary}

\begin{proof}
If the minimal ideal is all of $M$, the monoid is a finite group and
its positive cohomology at the indicated prime vanishes.  Otherwise
write its Rees matrix minimal ideal as $I\times H\times J$.
Collapse each minimal $H$-class, leaving all outside elements
unchanged, to obtain a monoid $\overline M$ with rectangular-band
minimal ideal.  The proof of
Lemma~\ref{lem:B-collapse-minimal-group} works over $k=\mathbb F_p$
whenever $|H|$ is invertible.  Explicitly, classwise averaging on
$kK\cong kMe\otimes_{kH}ekM$ is a bimodule idempotent.  Its kernel
$D$ is the kernel of $kM\to k\overline M$, is projective as a right
$kM$-module, and satisfies $D\otimes_{kM}k=0$.
The resolution $0\to D\to kM\to k\overline M\to0$ and
change of rings show that the actual quotient map
$BM\to B\overline M$ is a $k$-homology equivalence.  Its cohomology
isomorphism therefore preserves all Steenrod operations, not only
products.

The pointed Rees quotient $Q=\overline M/(I\times J)=M/K$ has
unit group $G$.  Its nonunits form an ideal $N$, and have no
nonzero idempotent: such an idempotent would lift to an idempotent
outside both regular classes of $M$.  A finite semigroup with zero
and no nonzero idempotent is nilpotent.  One direct proof is to take
a nonzero product longer than the number of nonzero elements.
Two of its nonzero prefixes coincide, say $a=at$ for a nonempty
intermediate word $t$.  Then $a=at^r$ for every $r$, so an
idempotent power of $t$ is nonzero, a contradiction.  Thus $N$ is
nilpotent, and Theorem~\ref{E:nilpotent-Rees-with-units} applies.
There is always a prime avoiding the two finite group orders,
whereas the suspended polynomial generator has an infinite
Steenrod orbit at every prime.
\end{proof}

\subsection{Nilpotent ideals with regular Rees quotients}

\begin{lemma}\label{E:regular-balanced-finiteness}
Let $B$ be a finite regular pointed monoid, with its zero as
basepoint, and put $R_B=\Sigma^\infty B$.  Let $p$ divide none of
the orders of its nonzero maximal subgroups.  If $L,V$ are a right
and a left $R_B$-module with finite-dimensional total
$\mathbb F_p$-homology, then
$L\wedge_{R_B}^{\mathbb L}V$ also has finite-dimensional total
$\mathbb F_p$-homology.
\end{lemma}

\begin{proof}
Induct on the number of nonzero $\mathcal J$-classes of $B$.
Choose a minimal nonzero class $J$ and an idempotent $e\in J$.
Since $B$ is regular, $J\cup\{0\}$ is a completely $0$-simple
principal ideal.  If $H$ is the maximal group at $e$, then
$eBe=H\cup\{0\}$ and multiplication gives a based bijection
\[
 Be\wedge_H eB\ \cong\ J\cup\{0\}.                 \tag{E.39}
\]
This is the Rees matrix normal form: after normalizing the sandwich
entry at $e$ to $1$, each element $(i,h,j)$ is uniquely a product
$(i,h,0)(0,1,j)$ modulo the intermediate $H$-coordinate.
The nonbasepoint parts of $Be$ and $eB$ are free on their indicated
$H$-sides.  Thus the derived balanced smash in (E.39) agrees with
the ordinary one after taking suspension spectra.

Write $R=R_B$, $K=\Sigma^\infty(J\cup\{0\})$, and
$A=\Sigma^\infty(B/J)$.  Then
\[
 K\simeq Re\wedge_{\mathbb S[H]}^{\mathbb L}eR,
 \qquad K\longrightarrow R\longrightarrow A,
 \qquad A\wedge_R^{\mathbb L}A\simeq A.               \tag{E.40}
\]
The last statement follows by tensoring the preceding cofiber
sequence, since $Ae=eA=0$.  Associativity gives the natural triangle
\[
 Le\wedge_{\mathbb S[H]}^{\mathbb L}eV
 \longrightarrow L\wedge_R^{\mathbb L}V
 \longrightarrow
 (L\wedge_R^{\mathbb L}A)
       \wedge_A^{\mathbb L}(A\wedge_R^{\mathbb L}V).   \tag{E.41}
\]
The idempotent summands $Le,eV$ have finite total mod-$p$ homology.
Extension to $H\mathbb F_p$ commutes with derived tensors and
homotopy orbits; the resulting modules have homology in a bounded
range, so the following group-homology spectral sequence converges.
The first term in (E.41) is the homotopy orbit spectrum of their
smash product under the finite group $H$.  The group-homology
spectral sequence collapses to coinvariants because $p\nmid|H|$;
its total mod-$p$ homology is therefore finite-dimensional.

The descended modules also retain finite total mod-$p$ homology.
For example, $L\wedge_R^{\mathbb L}A$ is the cofiber of
$Le\wedge_{\mathbb S[H]}^{\mathbb L}eR\to L$.
As a left $H$-set, $eB\setminus\{0\}$ is finite and free, so the
source is, on underlying spectra, a finite wedge of copies of $Le$.
The analogous assertion holds on the left.  The quotient $B/J$ is
again a finite regular pointed monoid with fewer nonzero classes,
and its nonzero maximal groups are among those of $B$.
Induction makes the last term in (E.41) finite in total mod-$p$
homology.  The long exact sequence proves the assertion for its
middle term.
\end{proof}

\begin{theorem}\label{E:nilpotent-ideal-regular-quotient}
Let $M$ be a finite monoid with rectangular-band minimal ideal $K$.
Suppose the pointed Rees quotient $Q=M/K$ has a nilpotent two-sided
ideal $N$ containing zero such that $Q/N$ is a regular pointed
monoid.  No section of $Q\to Q/N$ is assumed.
If $p$ divides none of the orders of the nonzero maximal subgroups
of $Q/N$, then $H^*(BM;\mathbb F_p)$ is locally finite over the
Steenrod algebra.  In particular $BM\not\simeq\Sigma\CP^\infty$.

For a minimal Rees matrix ideal over a finite group $H$, the same
conclusion holds provided also that $p\nmid|H|$.
\end{theorem}

\begin{proof}
First assume the minimal ideal is a band $I\times J$.  Use
Proposition~\ref{E:band-free-group-bar} to write
\[
 \widetilde{\Sigma^\infty BM}\simeq
 Z:=Y\wedge_A^{\mathbb L}X,
 \quad A=\Sigma^\infty Q,
 \quad X=\Sigma^\infty BF_I,
 \quad Y=\Sigma^\infty BF_J.                          \tag{E.42}
\]
We construct actual filtered ring and module spectra, rather than
only a filtration of their ordinary homology modules.

Choose $h\geq1$ such that $N^{h+1}=\{0\}$.  Give $A$ the finite
decreasing filtration
\[
 A^0=A,\qquad A^s=\Sigma^\infty N^s\quad(s\geq1),
 \qquad A^{h+1}=0.                                   \tag{E.43}
\]
Every $N^s$ is a pointed ideal, and products satisfy
$N^sN^t\subseteq N^{s+t}$.  Thus (E.43) is a filtered associative
ring spectrum, with the sphere unit in weight zero.
For the endpoints put
\[
 X^0=X,\qquad
 X^s=\Sigma^\infty
       \bigcup_{a\in N^s\setminus\{0\}}
              B\operatorname{im}\Phi_a\quad(s\geq1),  \tag{E.44}
\]
and define $Y^s$ in the same way for the right action.  An empty
union means the basepoint.  These are actual submodule spectra
under all of $Q$: composing any actor with an element of $N^s$
again gives an element of $N^s$, or zero.  More generally the
action maps factor strictly as
\[
 A^s\wedge X^t\longrightarrow X^{s+t},\qquad
 Y^t\wedge A^s\longrightarrow Y^{s+t},                 \tag{E.45}
\]
with the usual interpretation in weight zero.  These factorizations
respect all multiplication identities.  We therefore have coherent
filtered left and right modules, not merely maps whose homotopy
classes happen to satisfy (E.45).

All terms of these three filtrations are connective.  The endpoint
terms have finite CW homotopy type, by
Lemma~\ref{E:free-subset-intersections}; the ring terms are finite
wedges of spheres.  Their associated graded terms consequently
have finite-dimensional total mod-$p$ homology.

\emph{The filtered absolute bar.}
Take the ordinary derived bar over the sphere, with degree-$n$
term $Y\wedge A^{\wedge n}\wedge X$.  Use convolution of the
decreasing filtrations: the level-$L$ part is the homotopy colimit
of tensors of filtration terms whose indices sum to at least $L$.
For fixed $L$ this is a finite colimit calculation, since every
input filtration is zero above $h$; zero tails may equivalently
be truncated at $\max(h+1,L)$ in each index.
Equations (E.43)--(E.45) make all bar faces and degeneracies
filtered.  Let $F^LZ$ be the realization of its level-$L$ part.
Its level zero is $Z$ in (E.42).

In bar degree $n$ every nonzero filtration index is at most $h$,
and there are $n+2$ factors.  Thus the level-$L$ bar term is zero
whenever $L>(n+2)h$.  Every remaining term is connective, because
connective spectra are closed under smash products and colimits.
The skeletal realization filtration adds its degree-$n$ normalized
term with suspension $n$.  It follows that
\[
 F^LZ\text{ is at least }\lceil L/h\rceil-2
       \text{ connective}.                           \tag{E.46}
\]
In particular the connectivity tends to infinity with $L$.

\emph{Finite cohomology of each weight.}
For a decreasing filtration write
$\operatorname{gr}^s T=\operatorname{cofib}(T^{s+1}\to T^s)$.
The associated graded construction is strong monoidal for the
convolution just used:
\[
 \operatorname{gr}^w(U\wedge V)
 \simeq\bigvee_{s+t=w}
        \operatorname{gr}^sU\wedge\operatorname{gr}^tV. \tag{E.47}
\]
This follows by taking the successive cofibers of the finite
pushout diagrams defining convolution; smash products preserve
cofiber sequences.  Iteration gives the same formula for any
finite number of factors.  Associated graded also commutes with
geometric realization, since both cofibers and realization are
colimits.  Consequently $\operatorname{gr}Z$ is the bar over the
graded ring spectrum $\operatorname{gr}A$, with endpoints
$\operatorname{gr}Y,\operatorname{gr}X$.

The weight-zero ring is
\[
 R=\operatorname{gr}^0A
   =\operatorname{cofib}(\Sigma^\infty N\to\Sigma^\infty Q)
   \cong\Sigma^\infty(Q/N).                           \tag{E.48}
\]
In the \emph{graded} ring there is a canonical inclusion of this
weight-zero part.  Thus $\operatorname{gr}A=R\vee D$, where $D$
is supported in the strictly positive weights $1,\ldots,h$.
This splitting is in the associated graded ring; no multiplicative
section into $A$ is being asserted.

Compute the same graded derived tensor using the relative bar
over $R$.  Its normalized degree-$n$ term is
\[
 \operatorname{gr}Y\wedge_R^{\mathbb L}
      D\wedge_R^{\mathbb L}\cdots\wedge_R^{\mathbb L}D
      \wedge_R^{\mathbb L}\operatorname{gr}X,           \tag{E.49}
\]
with $n$ copies of $D$.  To see the normalization directly, expand
each internal copy of $R\vee D$.  Summands containing an $R$ are
exactly the degeneracy summands, and their latching cofiber is
the all-$D$ summand (E.49).

In a fixed weight $w$, (E.49) is a finite wedge of tensors
\[
 (\operatorname{gr}^{s_0}Y)\wedge_R^{\mathbb L}
 (\operatorname{gr}^{s_1}A)\wedge_R^{\mathbb L}\cdots
 \wedge_R^{\mathbb L}(\operatorname{gr}^{s_n}A)
 \wedge_R^{\mathbb L}(\operatorname{gr}^{s_{n+1}}X),
\]
where the indices sum to $w$, the two endpoint indices are
nonnegative, and all $n$ internal indices are positive.
Each factor has finite total mod-$p$ homology.  Since $Q/N$ is
regular, repeated application of
Lemma~\ref{E:regular-balanced-finiteness} gives the same finiteness
for each displayed tensor.  Moreover $n\leq w$, so only finitely
many normalized relative bar degrees occur in weight $w$.
Its finite skeletal filtration proves
\[
 \dim_{\mathbb F_p}H^*(\operatorname{gr}^wZ;\mathbb F_p)
       <\infty\qquad\text{for every }w\geq0.          \tag{E.50}
\]

For fixed $L$, the cofiber $Z/F^LZ$ has a finite filtration with
successive cofibers $\operatorname{gr}^wZ$, $0\leq w<L$.
Its total mod-$p$ cohomology is therefore finite-dimensional.
By (E.46), every class of $H^d(Z;\mathbb F_p)$, for fixed $d$,
is pulled back from one such cofiber when $L$ is sufficiently
large.  The cofiber map is an actual map of spectra.  Stable
Steenrod operations commute with it, so the entire orbit of the
class lies in the image of the finite-dimensional total
cohomology of this one cofiber.  This proves local finiteness.
The associated graded argument was used only for finite dimension
estimates; no identification of Steenrod operations on an
algebraic spectral sequence is required.

Equation (E.42) transfers the conclusion to $BM$.  If the minimal
ideal has a nontrivial group $H$, its actual classwise-averaging
quotient is a mod-$p$ homology equivalence for $p\nmid|H|$, as
proved in Corollary~\ref{E:two-regular-classes}.  It leaves $Q=M/K$
unchanged, so the band case applies.  Finally, a prime avoids all
the finitely many specified group orders, while the suspended
generator of $H^2(\CP^\infty;\mathbb F_p)$ has an infinite
Steenrod orbit at every prime.
\end{proof}

\begin{remark}\label{E:regular-quotient-scope}
This theorem includes the split regular-submonoid situation
$Q=B\cup N$, $B\cap N=\{0\}$, as well as the earlier finite-unit
case.  Its hypothesis still restricts the ideal structure.
A nonregular element may have a nonzero regular element below it
in the $\mathcal J$-order, preventing the nonregular elements from
lying in a nilpotent ideal.  Removing additional regular ideals
preserves finite underlying endpoint spectra, but does not by
itself supply the finite image filtrations (E.44) for their
coherent quotient actions.  The principal-ideal homotopy-colimit
construction of Lemma~\ref{E:regular-controlled-localization}
provides that additional step.  The resulting general theorem and
the finite-category reduction are proved in the following subsections.
\end{remark}


\subsection{A stable reduction from finite categories to finite monoids}

The following construction preserves the local groups and shifts the
entire reduced cohomology, including its stable operations.  It is a
reduction for the general local-finiteness question; it does not assert
an unstable equivalence between the resulting classifying spaces.

\begin{proposition}\label{E:category-consolidation}
Let $\mathcal C$ be a nonempty finite category with $v$ objects and $a$
arrows.  There is an explicitly defined monoid $M_{\mathcal C}$ of order
$1+a+(v+1)^2$ with rectangular-band minimal ideal and a natural stable
equivalence
\[
 \widetilde{\Sigma^\infty BM_{\mathcal C}}
       \simeq \Sigma^2\Sigma^\infty(B\mathcal C)_+ .
\]
Every nontrivial maximal subgroup of $M_{\mathcal C}$ is a maximal
subgroup of an endomorphism monoid of $\mathcal C$, and conversely.
In particular, $M_{\mathcal C}$ is aperiodic if and only if all
endomorphism monoids of $\mathcal C$ are aperiodic.
\end{proposition}

\begin{proof}
Write $O=\operatorname{Ob}\mathcal C$ and $I=J=O\amalg\{*\}$.
Let $Q$ be the consolidation of $\mathcal C$: its elements are the
arrows of $\mathcal C$, a zero $0$, and a fresh identity $1$.
Multiplication of arrows is categorical composition when defined and
zero otherwise.  For an arrow $f:x\to y$, define maps
\[
 \sigma_f(x)=y,\qquad \sigma_f(z)=*\quad(z\ne x),
 \qquad
 \tau_f(y)=x,\qquad \tau_f(z)=*\quad(z\ne y).
\]
Both maps fix $*$.  Set $\sigma_1=\tau_1=\mathrm{id}$ and let the
zero act constantly at $*$.  Then
$\sigma_{fg}=\sigma_f\sigma_g$ and
$\tau_{fg}=\tau_g\tau_f$.

Take the disjoint union
\[
 M_{\mathcal C}=\{1\}\amalg\operatorname{Mor}\mathcal C
                       \amalg(I\times J).
\]
Give $I\times J$ its rectangular-band multiplication
$(i,j)(i',j')=(i,j')$.  An arrow $f$ acts by
$f(i,j)=(\sigma_f(i),j)$ and
$(i,j)f=(i,\tau_f(j))$.  The product of two arrows is their
categorical composite when it exists and $(*,*)$ otherwise.
The fresh $1$ is the identity.  The displayed action identities,
together with the fact that every arrow fixes $*$ on both sides,
verify associativity, including products which enter the band.
The band is the minimal ideal: multiplying any band element on the
left and right by band elements reaches every band element, and any
nonempty ideal contains a band element.  The quotient by this ideal
is precisely $Q$.

Put $A=\Sigma^\infty Q$, with the zero interpreted as the basepoint.
The stable band-bar formula of
Proposition~\ref{E:band-free-group-bar} identifies the left-hand side
with the pointed two-sided bar on the row and column free-group
spaces.  Here every nonzero arrow fixes the distinguished element
$*$, so its free-group map sends each basis generator to a basis
generator or to $1$; it does not introduce a conjugating word.
The canonical map from the wedge of basis circles to the classifying
space of the free group is natural for precisely these maps.
Consequently the two endpoint $A$-modules in that formula are
respectively $\Sigma U$ and $\Sigma V$, where
\[
 U=\bigvee_{x\in O}\mathbb S,
 \qquad V=\bigvee_{x\in O}\mathbb S
\]
carry the right and left partial-permutation actions above.  Thus
\[
 \widetilde{\Sigma^\infty BM_{\mathcal C}}
       \simeq \Sigma^2(U\mathbin{\wedge_A^{\mathbb L}}V).
\]
This use of the free-group model supplies the coherent module
identification; no choice of an equivariant splitting of the original
row augmentation is being assumed.

For each object $x$, the categorical identity $e_x$ is a literal
idempotent of $Q$.  Distinct $e_x,e_y$ are orthogonal.  The full
spectral category on these idempotent summands has morphism spectra
\[
 e_y A e_x=\Sigma^\infty\mathcal C(x,y)_+,
\]
with the given categorical composition.  Both endpoint modules are
supported on these summands, and each $e_x$-summand is $\mathbb S$.
The standard idempotent-corner bar comparison therefore identifies
$U\wedge_A^{\mathbb L}V$ with the two-sided constant bar of this
spectral category.  To see the comparison directly, split the
orthogonal idempotents in the stable module category and insert their
representable modules between consecutive factors.  Spectral Yoneda
contracts each adjacent representable pair; the resulting simplicial
term in degree $n$ is
\[
 \bigvee_{x_0\xrightarrow{f_1}x_1\longrightarrow\cdots
                   \xrightarrow{f_n}x_n}\mathbb S.
\]
Its face and degeneracy maps are exactly those of the nerve.  The
complementary fresh-identity summand acts trivially on both endpoints
and contributes no tensor factor.  Realization gives
$\Sigma^\infty(B\mathcal C)_+$, as required.  In particular, the
formal sum of the $e_x$ is used only as a stable idempotent, not as a
set map or as an additional element of $Q$.

Every idempotent of $M_{\mathcal C}$ is either the fresh identity, a
band element, or an idempotent endomorphism of $\mathcal C$.  The first
has trivial unit group, as no product of two nonidentity elements is
the fresh identity.  Band elements have trivial local groups.  At a
categorical idempotent $e$, an invertible element of the corner cannot
belong to the band, since a product involving the band cannot equal
$e$.  Its units are therefore exactly the units of the categorical
corner $e\operatorname{End}_{\mathcal C}(x)e$.  This proves the
assertion about local groups.
\end{proof}

\begin{corollary}\label{E:monoid-category-LF-equivalence}
Fix a prime $p$.  The assertion that the mod-$p$ cohomology of every
finite monoid whose maximal subgroups have order prime to $p$ is
locally finite over the Steenrod algebra is equivalent to the same
assertion for all finite categories whose local subgroups have order
prime to $p$.
\end{corollary}

\begin{proof}
One implication is immediate by viewing a monoid as a one-object
category.  For the other, apply
Proposition~\ref{E:category-consolidation}.  Stable suspension shifts
degree and commutes with every stable Steenrod operation; it neither
creates nor removes an infinite-dimensional operation orbit.  The
additional degree-zero class of $(B\mathcal C)_+$ causes no problem.
\end{proof}

\begin{corollary}\label{E:category-regular-nilpotent-scope}
Let $\mathcal C$ be a finite category whose pointed consolidation $Q$
has a nilpotent ideal $N$ with regular quotient $Q/N$.  At every prime
avoiding the orders of the maximal subgroups of $Q/N$, the mod-$p$
cohomology of $B\mathcal C$ is locally finite over the Steenrod algebra.
In particular $B\mathcal C$ cannot have the homotopy type of
$\Sigma\mathbb{CP}^{\infty}$.
\end{corollary}

\begin{proof}
The minimal ideal of $M_{\mathcal C}$ is a band and its pointed Rees
quotient is $Q$.  Apply
Theorem~\ref{E:nilpotent-ideal-regular-quotient} and the stable
equivalence above.  There is always a prime outside the finite set of
local-group divisors, whereas the suspended integral projective-space
generator has an infinite Steenrod orbit at every prime.
\end{proof}

\subsection{Steenrod local finiteness at good primes}

Locally finite unstable modules are the first level of the usual
Krull filtration; see \cite[pp.~519--520]{schwartz2001krull}.
The argument below is given in full and does not apply the additional
hypotheses of that source's Theorem~0.1 to monoids.  The literature
search and exact attribution checks are recorded separately; no
priority claim is made.

All monoids in this subsection are finite pointed monoids: their
absorbing zero is the basepoint.  Actions and smash products are
pointed.  A finite group $H$ is treated as $H_+$.  Derived balanced
smash products are realized two-sided bars in pointed spaces, or
their suspension spectra when a stable comparison is specified.

\begin{definition}\label{E:controlled-action}
Fix $p$.  A pointed CW space with commuting cellular actions of finitely many
finite monoids is \emph{controlled} if every joint actor image is a
CW subcomplex and has finite total $\mathbb F_p$ homology.  The identity
actors are included, so the space itself has finite total homology.
We use the usual mapping-cylinder and simplicial-replacement CW
models; the explicit constructions below preserve this condition.
Only finite total homology, not a finite number of cells, is required.
\end{definition}

\begin{lemma}\label{E:LF-unstable-extension}
Locally finite unstable modules over the mod-$p$ Steenrod algebra
are closed under submodules, quotients, and extensions.  Consequently
a finite sequence of pointed homotopy pushouts and cofibers of spaces
with locally finite cohomology again has locally finite cohomology.
\end{lemma}

\begin{proof}
Submodules and quotients are immediate.  In an extension
$0\to K\to E\to Q\to0$, it suffices to consider a homogeneous
$x\in E$.  If $x$ lies in $K$ there is nothing to prove.  Otherwise take a finite homogeneous basis
of the Steenrod orbit of its image in $Q$.  Choose homogeneous lifts,
including $x$, of this basis.  Instability permits only finitely many
elementary squares, or powers and the Bockstein, on these finitely
many lifts.  Subtract the corresponding linear combinations of lifts
in their images.  The resulting finitely many errors lie in $K$.
Their orbits span a finite-dimensional submodule.  Its sum with the
span of the chosen lifts is finite-dimensional and closed under all
elementary operations, hence under the Steenrod algebra.  This
contains the orbit of $x$.

The cohomology exact sequence of a pointed cofiber or homotopy
pushout expresses the middle cohomology as an extension of
subquotients of the other cohomologies and their suspensions.
Suspension commutes with stable operations and preserves local
finiteness.  The middle cohomology is an unstable module, so the
extension assertion just proved applies.  No extension-closure
assertion for arbitrary stable modules is being used.
\end{proof}

\begin{lemma}\label{E:principal-image-replacement}
Let $R$ act on a controlled pointed space $X$ on the left, commuting
with an external finite-monoid action.  For a two-sided ideal
$P\subseteq R$, let $\mathcal P(P)$ be the finite poset of principal
right ideals $aR$ with $a\in P$, ordered by inclusion, and put
\[
 \mathcal F_P X=\mathop{\operatorname{hocolim}}_{aR\in\mathcal P(P)}
                         aX.
\]
The zero image is the one-point space; the homotopy colimit is
pointed.  Equal principal right ideals give literally equal images.
Then the following hold.
\begin{enumerate}
\item $\mathcal F_R X\to X$ is an $R$-equivariant underlying
      homotopy equivalence, also equivariant for the external action.
\item Ideal inclusions induce CW-subcomplex inclusions between the
      chosen models.  Left multiplication sends the diagram at $aR$
      to that at $baR$, strictly functorially.
\item Every joint actor image in $\mathcal F_P X$ has finite total
      mod-$p$ homology.  Its intersection with $\mathcal F_T X$, for
      a smaller ideal $T\subseteq P$, also has finite total homology.
      Consequently all joint images in
      $\mathcal F_P X/\mathcal F_T X$ have finite total homology.
\item If $P^h=0$, the spaces
      $F^sX=\mathcal F_{P^s}X$ ($P^0=R$) form a finite filtration,
      with strict multiplication maps
      $P^r\wedge F^sX\to F^{r+s}X$ for all $r,s\geq0$.
\end{enumerate}
There is a right-action version using principal left ideals.
\end{lemma}

\begin{proof}
If $aR=bR$, write $a=bu$ and $b=av$; then $aX=bX$.
If $aR\subseteq bR$, the corresponding map is the literal inclusion
$aX\subseteq bX$.  The object $1R$ is terminal in $\mathcal P(R)$,
proving the first assertion.  Multiplication $b$ sends $aR$ to
$baR$ and $aX$ to $baX$ and respects inclusions.  These maps compose
strictly.  External actors act objectwise and commute with them.

Here is a precise image description, including degeneracies.
The mapping-cylinder model is assembled from the finitely many
strict chains $a_0R\subsetneq\cdots\subsetneq a_nR$ and the pieces
$\Delta^n_+\wedge a_0X$.  A joint actor consisting of $b\in R$ and
an external actor $d$ sends such a chain to
$ba_0R\subseteq\cdots\subseteq ba_nR$.  Repeated vertices are
removed by the usual degeneracy map of simplices.  On the coordinate
space its image is all of
\[
 b d(a_0X)=d(ba_0X).
\]
If different chains have the same resulting chain, this coordinate
image depends only on its initial right ideal and on $d$, not on the
chosen lift.  Thus the actor image is the finite mapping-cylinder
realization over the simplicial image of the finite indexing nerve,
with these full coordinate-image spaces.  This simplicial image is
closed under faces.  A first face uses the inclusion of the initial
coordinate image into the next one; the other faces and the collapsed
degeneracies are precisely the ordinary diagram identifications.
Every coordinate image has finite total homology by the hypothesis
on joint actors.  A finite cellular realization therefore has finite
total homology.

Intersection with $\mathcal F_T X$ selects the finite subdiagram of
those image chains all of whose vertices belong to $\mathcal P(T)$.
The same reasoning applies.  These are CW-subcomplexes in the
mapping-cylinder model.  In the quotient the actor image is its
previous image divided by that intersection, so the cofiber exact
sequence proves the quotient assertion.  This also proves the claimed
CW-subcomplex condition for all the images constructed here.

For the final assertion, use the same construction also at $s=0$,
rather than replacing $F^0X$ by $X$.  The formula
\[
 b\cdot(a_0R\subsetneq\cdots\subsetneq a_nR,x)
   =(ba_0R\subseteq\cdots\subseteq ba_nR,bx)
\]
has its image in $F^{r+s}X$ when $b\in P^r$ and the original chain
lies in $P^s$.  Associativity and all filtration comparisons are
strict, since they are multiplication and inclusion of actual
ideals.  For $s\geq h$ the only image is the basepoint.  The right
version reverses multiplication.
\end{proof}

\begin{lemma}\label{E:regular-controlled-localization}
Let $J$ be a minimal nonzero regular $\mathcal J$-class ideal of a
finite pointed monoid $R$, and let $A=R/J$.  A controlled left
$R$-space $Y$, with a commuting external action, has a controlled
actual left $A$-space $Y'$ and an equivariant underlying equivalence
\[
 Y'\simeq A\mathbin{\wedge_R^{\mathbb L}}Y.
\]
The same holds on the right.  This comparison preserves the external
actions and can therefore be used inside longer balanced bars.
\end{lemma}

\begin{proof}
Choose $e\in J$ idempotent and let $H$ be its maximal group.
Rees coordinates give a literal pointed bimodule bijection
\[
 Re\wedge_H eR\ \xrightarrow{\ \cong\ }\ J.
\]
Both side sets are free on the relevant $H$-side, so this is also
their derived balanced product.  The evaluation map therefore gives
an actual mapping-cone model
\[
 Z=\operatorname{hocofib}(Re\wedge_H eY\longrightarrow Y)
       \simeq A\mathbin{\wedge_R^{\mathbb L}}Y.
\]
The bar preserves pointed homotopy cofibers; $eR\wedge_R^{\mathbb L}Y
\simeq eY$ follows by retracting the representable $R$-module.
The cone is an actual $R$-space.

Every joint actor image in $Z$ is the cone of the corresponding
image map.  Its source is a finite wedge of copies of a joint image
in $eY$: left multiplication on $Re$ has as its nonzero image a
union of free right $H$-orbits.  Its target is a joint image in $Y$.
Thus these cone images are CW-subcomplexes and have finite total
mod-$p$ homology.

For every nonzero $j\in J$, the image $jRe$ is one free right
$H$-orbit, and evaluation
\[
 jRe\wedge_H eY\longrightarrow jY
\]
is a homeomorphism.  Explicitly choose $u\in jRe$ and $v\in eR$
with $vu=e$ and $uv$ a left identity on the row containing $j$.
Then $u:eY\to jY$ and $v:jY\to eY$ are inverse homeomorphisms.
Consequently each pointed image $jZ$ is contractible: it is the cone
of this homeomorphism.  The zero image is also contractible.

Apply the principal-image construction to $Z$.  The subspace
$\mathcal F_JZ\subseteq\mathcal F_RZ$ is pointed contractible,
because every object of its pointed homotopy-colimit diagram is
contractible.  It is $R$-invariant, and every $j\in J$ sends all of
$\mathcal F_RZ$ into it literally.  Hence
\[
 Y'=\mathcal F_RZ/\mathcal F_JZ
\]
is an actual $A$-space.  The quotient map is an equivariant underlying
equivalence, because it collapses a contractible CW-subcomplex; and
$\mathcal F_RZ\to Z$ is an equivariant underlying equivalence.
The image description in Lemma~\ref{E:principal-image-replacement}
shows that $Y'$ is controlled, including all external actors.
This supplies strict quotient actions without assuming that arbitrary
coherent actions on finite spectra have controlled images.
\end{proof}

\begin{theorem}\label{E:general-controlled-multibar}
Let $R_1,\ldots,R_k$ be finite pointed monoids, and let
$X_0,\ldots,X_k$ be controlled pointed spaces with the commuting
left and right actions needed to form
\[
 Z=X_0\mathbin{\wedge_{R_1}^{\mathbb L}}X_1
       \mathbin{\wedge_{R_2}^{\mathbb L}}\cdots
       \mathbin{\wedge_{R_k}^{\mathbb L}}X_k.
\]
If $p$ divides none of the maximal-subgroup orders of these monoids,
then $\widetilde H^*(Z;\mathbb F_p)$ is locally finite over the
Steenrod algebra.
\end{theorem}

\begin{proof}
Induct on the finite multiset of the numbers of nonzero elements in
the monoids $R_i$.  Use the well-founded order in which one number
may be replaced by any finite list of strictly smaller numbers.
For example the ordinal sum $\sum_i\omega^{|R_i|-1}$, written in
decreasing order of exponents, strictly decreases under each such
replacement.  With no balancing monoids the assertion is the
finite-total-homology hypothesis on the one remaining space.
If a balancing monoid has $0=1$, unitality forces each adjacent
pointed module to be the one-point space, so its full bar is trivial.

Choose one balancing monoid $R$, and operate on its two adjacent
controlled bimodule spaces before performing the other bars.  Choose
a minimal nonzero $\mathcal J$-class ideal $J$ of $R$.

First suppose $J$ is regular.  Put $A=R/J$ and choose $e,H$ as in
Lemma~\ref{E:regular-controlled-localization}.  The ideal factorization
gives a pointed homotopy-cofiber sequence of balanced bars, with
terms at the chosen position
\[
 X e\mathbin{\wedge_H^{\mathbb L}}eY
   \longrightarrow X\mathbin{\wedge_R^{\mathbb L}}Y
   \longrightarrow X'\mathbin{\wedge_A^{\mathbb L}}Y'.
\]
Indeed $A\wedge_R^{\mathbb L}A\simeq A$: tensoring the cofiber
$J\to R\to A$ with $A$ kills $J$, since $eA=0$.
Apply the controlled strictification lemma to both adjacent modules
in the last term.  The remaining external actions commute with
these constructions.  The quotient monoid is smaller, and its local
groups are unchanged outside $J$, so the induction hypothesis applies
to the last term of the full multibar.

The first term can be absorbed into one new controlled bimodule.
For a finite group $H$ it is
\[
 EH_+\wedge_H(Xe\wedge eY),
\]
with the diagonal action using inversion on one side.  Since
$p\nmid|H|$, exact coinvariants show that this space has finite total
mod-$p$ homology.  Every external actor image is the same homotopy
orbit construction on the corresponding actual image subspaces;
it has finite total homology for the same reason and is a CW
subcomplex in the bar model.  Thus it is controlled.  This removes
the chosen balancing monoid entirely, and the induction hypothesis
applies to the remaining shorter chain.  The cofiber sequence and
Lemma~\ref{E:LF-unstable-extension} prove the assertion in the regular
case.

Now suppose $J$ is nonregular.  Then $J^2=0$, as follows directly
from finiteness.  The set $J^2$ is an $R$-ideal contained in the
minimal nonzero ideal $J$.  If it were nonzero, it would equal $J$,
and hence $J^n=J$ for every $n$.  Choose a nonzero product of more
than $|J|$ elements of $J$.  All its prefix products are nonzero,
so two prefix products agree, say $p_i=p_j=p_i u$ with $i<j$ and
$u\in J$.  A sufficiently large idempotent power $u^\omega$ then
satisfies $p_i u^\omega=p_i\ne0$.  It is a nonzero idempotent in
$J$, contrary to nonregularity.  Thus this is the square-zero ideal
case.  Filter the chosen monoid by
\[
 R^0=R,\qquad R^1=J,\qquad R^2=0.
\]
Replace the two adjacent endpoint spaces by their principal-image
models, using right ideals for the left action and left ideals for
the right action.  Their zeroth levels are equivariantly equivalent
to the original spaces.  Their first levels use the ideal $J$, and
their second levels are the basepoint.  All these inclusions are
cofibrations.  The associated-graded endpoint spaces are the actual
quotients $F^s/F^{s+1}$.  They are controlled for the external actions
and the quotient $A=R/J$, by
Lemma~\ref{E:principal-image-replacement}; the $J$ action is literally
zero on each quotient because it raises filtration.

Filter the absolute multibar by the total of these endpoint weights
and the weights of the letters of the chosen $R$-bar.  In degree $n$
of that bar, no weight exceeds $n+2$.  Hence the suspension spectrum
of its filtration part $F^LZ$ has connectivity tending to infinity
with $L$, even after realization in the other bar directions.  In
particular every fixed-degree cohomology class of $Z$ is pulled back
from an actual cofiber $Z/F^LZ$ for sufficiently large $L$.

It remains to prove that each such finite-weight cofiber has locally
finite cohomology.  Associated graded is strong monoidal for the
convolution filtration: the finite pushout-product cube identifies
the quotient in weight $w$ with the wedge over weight partitions of
the smash products of the corresponding quotients.  It also commutes
with realization.  These assertions are statements about actual
cofiber maps of filtered spaces, or equivalently their suspension
spectra, not about a homology spectral sequence preserving operations.

The graded algebra has weight-zero part $A$ and positive part $J$
in weight one, with $J^2=0$.  Its unit and augmentation split by
weight.  Relative bar normalization over $A$ therefore computes
weight $w$ by a finite simplicial realization: a nondegenerate term
has at most $w$ positive $J$ factors.  Every such term is a longer
chain of balanced bars in which the one occurrence of $R$ has been
replaced by finitely many occurrences of the strictly smaller monoid
$A$.  The intervening $J$ factors are finite pointed $A$-bisets;
the endpoint graded pieces are the controlled spaces just described;
all other factors are unchanged.  The induction hypothesis makes
the cohomology of every term locally finite.  Finite realization,
including its positive suspension shifts, preserves local finiteness
by Lemma~\ref{E:LF-unstable-extension}.  Thus every graded weight has
locally finite cohomology.  The finitely many weights below $L$ give
the same conclusion for $Z/F^LZ$.

Finally, an arbitrary class of $Z$ pulls back from one of these
finite-weight cofibers.  Pullback commutes with every Steenrod
operation, so its orbit is finite-dimensional.  This completes the
well-founded induction.
\end{proof}

\begin{corollary}\label{E:general-good-prime-LF}
Let $M$ be a finite monoid.  At every prime not dividing the order
of any maximal subgroup of $M$, the mod-$p$ cohomology of $BM$ is
locally finite over the Steenrod algebra.  The same statement holds
for finite categories and their local subgroups.
\end{corollary}

\begin{proof}
Adjoin a new absorbing basepoint to $M$, obtaining the pointed monoid
$M_+$.  Let both endpoints be $S^0$, with all elements of $M$ acting
as the identity and the basepoint acting as zero.  These are finite
controlled spaces, and their pointed two-sided bar is $(BM)_+$.
Apply Theorem~\ref{E:general-controlled-multibar}.  The categorical
assertion follows from the stable consolidation construction in
Proposition~\ref{E:category-consolidation}, which preserves local
groups and shifts the entire cohomology with its stable operations.
\end{proof}

\begin{corollary}\label{E:Sigma-CP-infty-excluded}
The space $\Sigma\mathbb{CP}^{\infty}$ is not the classifying-space
homotopy type of a finite monoid or a finite category.
\end{corollary}

\begin{proof}
There are only finitely many primes dividing the local subgroup
orders of a finite category.  Choose any other prime $p$.  If
$u\in H^2(\mathbb{CP}^{\infty};\mathbb F_p)$ is the usual generator,
then the suspended class $\sigma u$ has an infinite Steenrod orbit.
At $p=2$ the successive operations
$\mathrm{Sq}^{2},\mathrm{Sq}^{4},\mathrm{Sq}^{8},\ldots$ give
$\sigma u^2,\sigma u^4,\sigma u^8,\ldots$.  At an odd prime the
successive powers $\mathcal P^1,\mathcal P^p,\mathcal P^{p^2},\ldots$
give $\sigma u^p,\sigma u^{p^2},\sigma u^{p^3},\ldots$.
They are nonzero in distinct degrees, contradicting
Corollary~\ref{E:general-good-prime-LF}.
\end{proof}

\begin{remark}\label{E:general-audit}
The complete proof above was independently checked by the root and a
second proof agent, including the strict image construction, the
pointed-space localization triangle, and the fixed-weight multibar
induction.  The finite cellular control in
\path{code/a_general_good_prime_control.py} checks the case where the
ordinary union of two contractible regular-ideal images is a circle,
while their pointed principal-ideal homotopy colimit is contractible.
It verifies every action composition, chain-map equation, and strict
quotient action in that example.  This is a corroborating finite
control, not the source of the all-order theorem.  The proof makes
no claim of literature priority and has no census-completeness
dependency.
\end{remark}

